\documentclass[a4paper]{article}
\usepackage{titlesec}
\usepackage{geometry}
\geometry{hmargin={2cm,2cm}, vmargin={3.5cm,3.5cm}}
\usepackage{graphicx}
\usepackage[table]{xcolor}
\usepackage{hyperref}
\definecolor{logo}{RGB}{20,99,163}
\newcommand{\sectionbreak}{\clearpage}
\usepackage[T1]{fontenc}
\usepackage[utf8]{inputenc}
\usepackage[english,italian]{babel}
\usepackage{fancyhdr}
\usepackage{lastpage}

% creazione comando subsubsubsection
\titleclass{\subsubsubsection}{straight}[\subsection]

\newcounter{subsubsubsection}[subsubsection]
\renewcommand\thesubsubsubsection{\thesubsubsection.\arabic{subsubsubsection}}
\titleformat{\subsubsubsection}
  {\normalfont\normalsize\bfseries}{\thesubsubsubsection}{1em}{}
    \titlespacing*{\subsubsubsection}
    {0pt}{3.25ex plus 1ex minus .2ex}{1.5ex plus .2ex}

\makeatletter
\def\toclevel@subsubsubsection{4}
\def\l@subsubsubsection{\@dottedtocline{4}{7em}{4em}}
\makeatother

\setcounter{secnumdepth}{4}
\setcounter{tocdepth}{4}
%fine creazione comando


%Comandi di impaginazione uguale per tutti i documenti
\pagestyle{fancy}
\lhead{\includegraphics[scale=0.2]{../../../../Images/logo}}
%Titolo del documento
\rhead{\large Verbale Interno del 2021-03-16}
%\rfoot{\thepage}
\cfoot{Pagina \thepage{} di \pageref{LastPage}}
\setlength{\headheight}{35pt}
\renewcommand{\footrulewidth}{0.4pt}




\begin{document}
\begin{titlepage}
    \begin{center}
        \includegraphics{../../../../Images/logo}\\
        \vspace{20px}
        \textcolor{logo}{\hrulefill}\\
        \vspace{20px}
        \textbf{\huge\textcolor{logo}{Verbale Interno del 2021-03-16}}\\
        \vspace{10px}
        \textcolor{logo}{\hrulefill}\\
        \vspace{40px}
        \textbf{\Large Informazioni sul documento}\\
        \vspace{20px}
        \begin{tabular}{p{100px} | p{100px}}
            \textbf{Versione} & 1.0.0\\
            \textbf{Responsabile} & Samuele De Simone\\
            \textbf{Redattore} & Samuele De Simone\\
            \textbf{Verificatore} & Nicolò Giaccone\\
            \textbf{Destinatari} & Prof. Tullio Vardanega \newline Prof. Riccardo Cardin, \newline TechSWEave\\
            \textbf{Stato} & Approvato\\
            \textbf{Uso} & Esterno\\
        \end{tabular}\\
        \vspace{60px}
        \href{mailto:techsweave@gmail.com}{techsweave@gmail.com}\\

    \end{center}
    \end{titlepage}

    %DIARIO DELLE MODIFICHE
    \begin{center}
        \textbf{\Large Diario delle modifiche}\\
        \vspace{10px}
        \begin{table}[h!]
        \centering
        \rowcolors{2}{logo!10}{logo!40}
        \renewcommand{\arraystretch}{1.8}
        \begin{tabular}{p{150px} p{90px} p{60px} p{60px} p{45px}}
            \rowcolor{logo!70} \textbf{Modifica} & \textbf{Autore} & \textbf{Ruolo} & \textbf{Data} & \textbf{Versione}\\
            Stesura Verbale & Samuele De Simone & Responsabile & 2021-03-17 & 0.0.1 \\    
            Completamento verbale & Samuele De Simone & Responsabile & 2021-03-17 & 0.0.2\\
            Verifica verbale & Nicolò Giaccone & Verificatore & 2021-03-18 & 0.1.0\\
            Approvato & Samuele De Simone & Responsabile & 2021-03-26 & 1.0.0\\
        \end{tabular}
    \end{table}
    \end{center}

    %INDICE
    \newpage
    \tableofcontents
    \newpage
    % CORPO DEL DOCUMENTO
    \section{Informazioni riunione}
    \begin{itemize}
        \item Data e ora: 2021/03/16 ore 18.00-19.00;
        \item Luogo: Canale Discord del gruppo tramite videochiamata;
        \item Partecipanti interni: Samuele De Simone, Manuel Varo, Tito Scutari, Nicolò Giaccone, Amedeo Meggiolaro, Marco Barbaresco;
        \item Partecipanti esterni: nessuno;
        \item Segretario: Amedeo Meggiolaro.
    \end{itemize}
    \section{Argomenti della riunione}
    \begin{itemize}
        \item Discussione sugli argomenti della \textit{Analisi dei Requisiti};
        \item Scelta dei casi d'uso da proporre nella riunione con i \textit{Proponenti};
        \item Discussione su come il gruppo si dividerà il lavoro per la stesura dell'Analisi dei Requisiti e dei Piani di Qualifica e di Progetto;
        \item Discussione sulle richieste da porre ai proponenti e sui dubbi da porre.
    \end{itemize}
    \section{Resoconto}
    
        \subsection{Discussione sugli argomenti della Analisi dei Requisiti}
        Il gruppo per prima cosa ha discusso sugli argomenti da inserire nel documento relativo all'analisi dei requisiti.
        \subsection{Scelta dei casi d'uso da proporre nella riunione con i proponenti}
        Il gruppo successivamente ha analizzato degli e-commerce online per trovare degli spunti per qualche caso d'uso da inserire nel prodotto, da proporre ai proponenti.
        \subsection{Divisione del lavoro per la stesura dei documenti}
        È stato deciso come suddividere tra i vari membri il lavoro di stesura dei documenti, in base alla loro importanza.
        \subsection{Discussione sulle richieste da porre ai proponenti e sui dubbi da porre}
        Alcuni membri del gruppo hanno esposto a tutti le loro perplessità in modo da discutere su che domande fare ai proponenti.
        
\end{document}
